\section{Results}
\label{sec:results}

% PROVISIONAL: the global four-way comparison (Sec. 4.2) is currently N=28; the
% 10 complementary reservoirs are being processed and will raise it to N=38.
% Numbers flagged [N=28] update on their arrival; the A/P curve (Sec. 4.1) is
% already N=42 (28 v4 + 14 v3, dual vs JRC). Figures are pulled live from
% analysis/method_comparison_output/ via \graphicspath.

\subsection{A/P predicts SAR monitorability a priori}
\label{sec:res_ap}

Across the pooled sample (42 reservoir--reference series: the global pilot plus
the four Sicilian reservoirs, all scored against their optical reference), the
static area-to-perimeter ratio is a strong predictor of surface-area skill:
Spearman $\rho(\text{A/P},\text{KGE}) = 0.57$ ($p = 10^{-4}$;
Figure~\ref{fig:ap_curve}). The relationship behaves as a \emph{ceiling}
rather than a simple trend. High-A/P reservoirs ($\geq$250~m) are reliably
tracked (median KGE 0.88, and three-quarters reach KGE~$\geq$~0.5),
whereas low-A/P reservoirs ($<$120~m) show both a lower median (0.49) and a
much wider spread (down to negative KGE), with only half reaching
KGE~$\geq$~0.5. This is the signature expected if A/P sets the geometric
ceiling on monitorability: compact shorelines leave little room for
mixed-pixel error, so skill is high and stable, while elongated, dendritic
shorelines merely \emph{permit} large errors that a second, radiometric axis
(Section~\ref{sec:res_wind}) then determines. Two high-A/P outliers with low
KGE (Rockport, Boegoeberg) are consistent with this reading: their geometry is
favourable, but ice cover and extreme drawdown, respectively, drive their
errors, failures that A/P cannot see.

\begin{figure}[t]
  \centering
  \includegraphics[width=0.9\linewidth]{ap_kge_curve_pooled}
  \caption{Surface-area skill (KGE of the dual-polarisation SVM against the
  optical reference) versus static A/P for the pooled sample ($n=42$).
  A/P sets a monitorability ceiling: high-A/P reservoirs are reliably tracked,
  low-A/P reservoirs are permitted, but not guaranteed, large errors.}
  \label{fig:ap_curve}
\end{figure}

\subsection{Adaptivity, not polarisation count, is the decisive lever}
\label{sec:res_fourway}

Comparing the four detectors under the identical pipeline overturns the
intuitive ``dual-pol beats single-pol'' expectation. Ranked by median KGE
across the global pilot [N=28], the order is
\emph{per-scene SVM} (0.64) $>$ \emph{per-scene VV Otsu} (0.56) $>$
\emph{VV Otsu without vectorisation} (0.52) $>$ \emph{fixed-training dual SVM}
(0.50) (Figure~\ref{fig:fourway}). The dual-polarisation learner with a fixed
2023 baseline ranks \emph{last} (below the single-polarisation threshold)
because the VV Otsu adapts to every scene by construction while the fixed SVM
does not. The decisive axis is therefore per-scene adaptivity, not the number
of polarisations.

Two like-for-like contrasts isolate the two axes. \emph{(i) Adaptivity},
holding polarisation fixed at dual: retraining the SVM per scene beats the
fixed 2023 baseline by a median $\Delta$KGE of $+0.066$ (16 win, 6 tie, 6 lose
of 28), so the arbitrary baseline year is at worst neutral and on median a
gain, which disposes of the ``why 2023?'' objection. \emph{(ii)
Polarisation}, holding adaptivity fixed (per-scene SVM vs per-scene Otsu):
overall the two are statistically indistinguishable (median $\Delta$KGE
$+0.01$), \emph{but the dual-pol advantage is significantly concentrated at low
A/P}: Spearman$(\text{A/P},\,\text{KGE}_{\text{dual}}-\text{KGE}_{\text{VV}}) =
-0.41$ ($p = 0.03$), with a low-A/P median gain of $+0.03$ collapsing to $0.00$
at high A/P. VH earns its keep only where mixed-pixel, low-water conditions
defeat a single-band threshold; elsewhere the extra polarisation is redundant.

\begin{figure}[t]
  \centering
  \includegraphics[width=\linewidth]{kge_4way}
  \caption{Four-detector KGE across the global pilot [N=28, updating to 38].
  Per-scene adaptivity (Otsu; adaptive SVM) outranks the fixed-training dual
  SVM; the dual-pol edge over VV concentrates at low A/P.}
  \label{fig:fourway}
\end{figure}

\subsection{Independent confirmation against PlanetScope near-truth}
\label{sec:res_sicily}

Validated against 3~m PlanetScope near-truth (which breaks the 30~m
JRC inter-product circularity), the four Sicilian reservoirs reproduce every
global finding (Figure~\ref{fig:sicily}). Mean KGE ranks per-scene SVM (0.80)
$\approx$ fixed dual (0.80) $>$ Otsu (0.71) $>$ no-vectorisation (0.65). The
adaptive and fixed dual SVMs tie against truth (mean $\Delta$KGE $+0.004$),
again showing the baseline year is harmless. Most tellingly, the dual-pol gain
over Otsu is largest at the lowest-A/P reservoir, Ancipa ($\Delta$KGE $+0.33$;
Otsu $0.23 \rightarrow$ adaptive $0.55$): there the single-band Otsu
\emph{floors} near 60~ha while the true surface falls to $\sim$5~ha, and the
per-scene dual SVM follows the drawdown down, direct visual evidence that
VH rescues the mixed-pixel, low-water regime that A/P flags as hardest.

\begin{figure}[t]
  \centering
  \includegraphics[width=\linewidth]{sicily_planet_validation}
  \caption{Surface-area time series of the four detectors against PlanetScope
  3~m near-truth for the Sicilian reservoirs. At low-A/P Ancipa the VV Otsu
  floors at low water while the per-scene dual SVM tracks the drawdown.}
  \label{fig:sicily}
\end{figure}

\subsection{Computational cost trades off against accuracy along A/P}
\label{sec:res_cost}

An audit of Earth Engine compute usage (EECU-seconds per task, paired by
reservoir, cache-confounded pairs removed) shows that cost is driven by
per-scene retraining, not by the number of polarisations
(Figure~\ref{fig:eecu}). Polarisation count alone is cheap: a dual SVM trained
once on a fixed year costs a median $0.69\times$ the single-band VV Otsu, whose
per-scene 256-bin histogram and threshold search outweigh a fixed classifier.
Making the dual classifier accurate, however, requires re-training it on every
scene, and this per-scene dual SVM is the most expensive method, at
$1.21\times$ the VV Otsu and $1.64\times$ the fixed-training dual. The full cost
ordering, fixed dual ($0.69\times$) $<$ fast ($0.77\times$) $<$ VV Otsu
($1.00\times$) $<$ per-scene dual ($1.21\times$), is close to the inverse of the
accuracy ordering (Section~\ref{sec:res_fourway}). The single-band ``fast'' path
(no vectorisation) is the cheapest option that still keeps a classifier, but its
$\sim$25\% saving over VV Otsu is not free: it degrades accuracy by a median
$\Delta$KGE of $-0.07$ globally and $-0.05$ against PlanetScope
(Sections~\ref{sec:res_fourway} and~\ref{sec:res_sicily}). Because the per-scene
dual SVM is at once the most accurate and the most expensive, and its accuracy
advantage is confined to low A/P (Section~\ref{sec:res_fourway}), the
cost-optimal choice depends on A/P: the per-scene dual SVM earns its premium
only where shorelines are complex, while the per-scene VV Otsu delivers equal
accuracy at lower cost elsewhere.

\begin{figure}[t]
  \centering
  \includegraphics[width=0.8\linewidth]{eecu_svm_vs_vv}
  \caption{Per-reservoir Earth Engine compute cost of the fixed-training dual
  SVM versus VV Otsu (log--log). Most points fall below the 1:1 line, so
  polarisation count alone is not the cost driver; the per-scene retraining that
  makes the dual classifier accurate, by contrast, makes it the most expensive
  method (Section~\ref{sec:res_cost}).}
  \label{fig:eecu}
\end{figure}

\subsection{Wind: the radiometric axis does not covary with geometry}
\label{sec:res_wind}

Finally, we test whether wind (the main radiometric failure mode) can be
folded into the geometric picture, and find it cannot. The Bragg hypothesis
predicts that wind roughening makes VV over-estimate land (under-detect water)
relative to dual-pol; across the pilot this correlation has the predicted
negative sign in \emph{none} of the reservoirs, and where significant it runs
the opposite way (Figure~\ref{fig:wind}). Nor does wind explain the
between-method accuracy gap: $r(\text{wind}_{p90}, \Delta\text{KGE})$ is not
significant, and the only significant predictor of $\Delta$KGE remains A/P
(Section~\ref{sec:res_fourway}). The wind and A/P axes are, as argued in
Section~\ref{sec:introduction}, genuinely orthogonal: A/P is a clean,
computable-a-priori predictor of the geometric ceiling, whereas the
radiometric residual (wind, ice, vegetation) is real but does not organise
along geometry and is not captured by choosing one polarisation over another.

\begin{figure}[t]
  \centering
  \includegraphics[width=0.9\linewidth]{delta_kge_vs_wind}
  \caption{Between-method skill difference versus wind exposure. Wind does not
  predict $\Delta$KGE; the geometric predictor A/P does. The radiometric axis
  is orthogonal to geometry.}
  \label{fig:wind}
\end{figure}
